\chapter{绪论}
%
\section{研究背景和意义}
\subsection{研究背景和意义}
%
关于\LaTeX{}以及基于\LaTeX{}写作的好处不再赘述。\LaTeX{}的入门资料推荐文献\parencite{_g}以及文献\parencite{_c}。

这里主要是想推荐一种“学术生态”，即利用各种工具展开科研工作，以达到事半功倍的效果。需要用到以下软件：
\begin{enumerate}[topsep = 0 pt, itemsep= 0 pt, parsep=0pt, partopsep=0pt, leftmargin=44pt, itemindent=0pt, labelsep=6pt, label=(\arabic*)]
	\item 	参考文献管理软件zotero\cite{_m}。很多人使用过endnote，但其实zotero也非常强大，强烈推荐。可到b站观看Struggle with Me出品的视频教程\cite{_k}入门（或其他最新教程，刚开始不推荐使用插件，会增加学习难度）。zotero自带pdf阅读器，也可以设置为使用其他阅读器。在zotero可以打开文件所在位置，故不推荐更改zotero的文件系统（尤其不推荐使用zotfile插件，事实上各种五花八门的插件增加了复杂性，实际上没有带来太多便利性）。理论上只需要包含文献元数据信息的bib文件（可以手动一篇一篇文章地收集）即可使用此模板，因此模板不依赖于任何参考文献管理软件，endnote用户或不使用参考文献管理软件的用户可以忽略本文zotero部分的讲解。
	\item	可截图获取文献中公式的软件mathpix\cite{_h}。在阅读别人的论文时，很可能需要把文章中的公式抄下来放到自己的笔记中，方便以后组会报告甚至论文中使用，这时使用mathpix可直接截图获取\LaTeX{}源码，非常方便。随着mathpix的使用成本越来越高，免费次数越来越少，2023起已经不再推荐。目前开源/免费的替代工具为：\href{https://www.simpletex.cn/}{SimpleTex}和\href{https://p2t.breezedeus.com/}{Pix2Tex}。现在（2025年）SimpleTex网页不收费但客户端早已经收费，Pix2Tex虽然一般般但是好过没有，还是不错的选择，毕竟一般来说公式不会太复杂。Pix2Tex也一直在改进，可以本地部署也可以在线使用，MacOS还有桌面应用程序。
	\item	TeXlive202x、TeXstudio（2022起推荐vscode），相当于开发环境和IDE。本模板是基于TeX的发行版TeXlive202x和编辑器TeXstudio进行的，百度这两个关键字分别安装。关于TeXstudio的使用（快捷键等）可另行查找资料。模板还支持更多ide，更多编译方式见GitHub首页readme.md。若在其他窗口打开了编译生成的pdf文件，记得关掉再编译，否则报错。TeXstudio的设置见第二章。vscode的设置见github讨论区的\href{https://github.com/mengchaoheng/SCUT_thesis/discussions/6}{vscode配置}。
\end{enumerate}

本文的章节安排如下：

第一章，绪论。

第二章，模板简介。主要介绍各文件的内容。

第三章，常用环境。介绍论文写作中常用的环境，包括：图、表、公式、定理。基本涵盖了常用的命令。

%第三章，参考文献设置。本模板对旧版的改动主要是参考文献部分，本章将简单参考文献设置以及
%编译选项的设置等等。

\section{国内外研究现状}
目前针对高频电刀的研究主要集中在以下几个方面：

一是对高频电刀的能量发生结构电路进行研究：如伊利诺伊大学芝加哥分校Congbo Bao等人使用GaN功率开关器件设计了一种高频交流发生装置，包含全桥高频逆变器和多谐振频率滤波器（MRF），通过隔离升压变压器作用在负载上，并通过一个闭环PI-前馈控制器进行输出电压控制，所提出的电路结构具有极低的总谐波失真（THD）和宽负载适应能力\cite{baoMultiresonantFrequencyFilterElectrosurgery2022,baoGaNHEMTBasedVeryHighFrequency2021}；赵忠源等人提出了一种以对称式E类射频放大器为核心的ESU结构，通过DAC控制直流电源输出从而控制设备的输出功率，并使用PID控制算法对输出功率进行控制\cite{ZhaoZhongYuanYongYuDanJiQieGeDeZhiNengDianWaiKeSheBeiDeYanZhi2022}；Sudip Mazumder等人汇总了现有文献中已知并公认的高频逆变方案：如E类、EF$_n$类、$\Phi_{2}$、 LCLC和双轨多相Buck，经过分析得出上述方案具有负载范围窄、开关损耗高和直接输出方波等缺陷，从而引出基于全桥逆变和多谐振频率滤波器的高频逆变器设计\cite{mazumderElectrosurgeryPowerElectronics2023}；Liu Liu等人提出了一种多相交错可重构的高频逆变器，电路结构可实现全桥和半桥切换从而实现宽范围电压输出，可提升动态性能、降低输出电压纹波与ZVS输出，但此结构下的输出电压和电流均为方波波形\cite{liuMultiphaseInterleavedReconfigurable2022}；

二是




